\clearpage
\section{The Propers: Themes and Movement}
\zlabel{triptych:concise:themes:start}
God gives his vineyard life and asks it to bear justice; the wounded vineyard asks him for the life it cannot give itself. The Sunday moves between those two directions. One reading follows the priority of gift into fruitful charity. Another follows the plea for visitation into a people renewed in Christ. They meet in the Son whom the stewards reject and in the nourishment that forms his members.

\subsection{A gift precedes the demand}
The Entrance's Mordecai speaks for endangered Israel. Creation belongs to the Lord, while the people cannot secure their own survival. The Collect approaches divine generosity from another need: forgiveness and gifts beyond the petitioner's merit or reach. Dependence is therefore already present before Isaiah's accusation. Neither the threatened nor the guilty must first become self-sufficient in order to pray.

Isaiah's vineyard has received loving preparation. Its failure is measured against that care, and its fruit has a definite name: justice. The outcry in place of justice exposes more than disappointing productivity. Chrysostom, explaining Matthew's related image, observes that the householder prepares and equips the vineyard before asking the tenants for anything. His gifts make their refusal more culpable. Augustine reaches the same order from John: Christ chooses disciples so that they may become good and bear fruit, not because their prior fruit has bought his choice.\footnote{Chrysostom, \work{Homilies on Matthew} 68.1; Augustine, \work{Tractates on John} 86.2--3.}

\subsection{The accused vineyard becomes a praying people}
The psalm changes the voice. Those who have heard the owner's accusation now speak from within the broken enclosure. They remember the vine brought from Egypt and ask its planter to look upon it and give it life. The reading and response hold judgment beside lament. Their conjunction does not make every wounded person's suffering evidence of equal guilt.

For the reading centered on gift, this plea shows why response remains dependent: the people invoke God because he quickens them. For the reading centered on restoration, it preserves the identity of the planting through disaster. Augustine and Bellarmine find its perfection in Christ. Bellarmine names Jewish apostles and converts, together with grafted Gentiles, as the concrete shape of renewal. The vineyard is reformed, not discarded for an unrelated people.\footnote{Augustine, \work{Exposition of Psalm} 79, §§9--11 (NPNF heading Psalm 80); Bellarmine, \work{Commentary on the Book of Psalms}, Psalm 79, at vv. 14--18.}

Matthew makes the crisis sharper. The tenants seek possession by killing the heir, but the rejected stone becomes the cornerstone. Rejection cannot determine the Son's place in the Father's work. The transfer of stewardship retains the demand for fruit; the Christian hearer acquires no immunity from it. Chrysostom distinguishes Isaiah's accusation of the vineyard from Jesus' accusation of its rulers. Both injustice and the stewards' resistance to correction matter, without making every member of the people a violent tenant.

\clearpage
\subsection{Thanksgiving becomes practiced peace}
Philippians follows its own apostolic course; Isaiah is the reading selected in relation to the Gospel. Paul's prayer, thanksgiving and attention to worthy things nevertheless give a concrete form to the response the vineyard owes. His final imperative is to practice what the community has learned, received, heard and seen. Thoughts must become deeds. Chrysostom explains peace through God's reconciliation of enemies in the gift of the Son, so peace begins deeper than a managed feeling.\footnote{\work{General Introduction to the Lectionary}, 66--68,106--107; Chrysostom, \work{Homilies on Philippians} 14.}

The two readings place different weight on Paul's exhortation. Thanksgiving acknowledges received gifts and frees obedience from the effort to purchase favor. Reconciliation requires a damaged community to live truthfully with its members. The preceding appeal to Euodia and Syntyche belongs to the Epistle's context: peace concerns actual relationships. Truth and justice remain among the things Paul commends; peace cannot depend upon concealing harm.

The adapted Johannine acclamation directs both movements toward lasting fruit. Augustine reads the following command of mutual love as its explanation. Fruit is charity, not the acquisition of an impressive religious reputation. Christ's choosing therefore calls forth action while removing superiority toward the neighbor. The Mass's common prayers belong to all three cycles; their generosity can nourish this response without turning Paul's course into another vineyard parable.

\subsection{Nourishment reaches beyond the moment of reception}
The Prayer over the Offerings joins the Church's obedience to redemption's sanctifying work. Aquinas explains the Eucharist's efficacy from Christ and his Passion: the sacrament sustains life and arouses charity to act. The Prayer after Communion seeks a life shaped by what has been received. Gift and response thus remain together at the altar. Worship does not make God a debtor; reception does not release the recipient from love.\footnote{Aquinas, \work{Summa theologiae} III, q. 79, a. 1, body and replies 2--3. The offerings analysis rests on the bounded Latin witness; its exact approved English remains uncollated.}

The Communion antiphons offer two accents. Lamentations speaks of seeking the good Lord amid a chapter of grief. Fruit may remain hidden, and repair may still be awaited. The Corinthians alternative names common participation and the one body. Chrysostom asks why people nourished by the same Christ fail to show the same love. Patience and incorporation are compatible, but these are alternatives; neither has been selected for an actual celebration.

Both directions end beyond immediate success. Augustine makes God himself the reward; Bellarmine distinguishes present grace from future glory. Aquinas joins present bodily service to future incorruption. Fruitfulness is therefore more than visible achievement, and restoration more than recovered strength. The people who receive Christ remain called to justice while awaiting a fulfillment they cannot manufacture.
\zlabel{triptych:concise:themes:end}
